% !TEX root = ../main.tex
\chapter{Análisis Asintótico y Secciones Matriciales}\label{sec:matricial}

\section{Matriz de Hessenberg y secciones truncadas}

\begin{defn}\label{def:hessenberg}
	Sea $(\Poly, \langle \cdot,\cdot\rangle_S)$ y sea $\{\phi_n\}_{n\ge0}$ la sucesión de polinomios ortonormales asociada al producto de Sobolev. La representación matricial de $\D$ es
	\[
		\mathbf D=(d_{i,j})_{i,j\ge0},
		\qquad
		d_{i,j}=\langle z\phi_j,\phi_i\rangle_S.
	\]
\end{defn}

\begin{prop}\label{prop:hessenberg}
	Sea $(\Poly, \langle \cdot,\cdot\rangle_S)$ y sea $\{\phi_n\}_{n\ge0}$ la sucesión de polinomios ortonormales asociada al producto de Sobolev.
	Entonces la matriz $\mathbf D$ es de Hessenberg, es decir,
	\[
		d_{i,j}=0,
		\qquad i>j+1.
	\]
\end{prop}

\begin{proof}
	La expansión de $z\phi_j$ en la base ortonormal solo puede involucrar polinomios de grado a lo sumo $j+1$. Por ortogonalidad, sus coeficientes frente a $\phi_i$ se anulan cuando $i>j+1$.
\end{proof}

\begin{defn}\label{def:seccion_truncada}
	Para cada $n\in\mathbb N$, definimos la \emph{sección principal} de orden $n$ de $\mathbf D$ por
	\[
		\mathbf D_n:=(d_{i,j})_{0\le i,j\le n}\in\mathcal M_{n+1}(\mathbb C).
	\]
	Equivale a la compresión proyectada de $\D$ al subespacio
	\[
		\mathcal P_n:=\operatorname{span}\{\phi_0,\dots,\phi_n\},
	\]
	es decir,
	\[
		\mathbf D_n\ \widehat{=}\ P_n\D|_{\mathcal P_n},
	\]
	donde $P_n$ denota la proyección ortogonal sobre $\mathcal P_n$.
\end{defn}

%%La importancia de $\mathbf D_n$ es doble. Por un lado, conserva la estructura de Hessenberg de $\mathbf D$ y por tanto retiene la información esencial del operador en dimensión finita. Por otro, ya no actúa sobre un espacio abstracto sino sobre $\mathbb C^{n+1}$, de modo que sus valores singulares y autovalores son objetos directamente calculables.

\begin{prop}\label{prop:charpoly_truncada}
Sea $\{\phi_n\}_{n\ge0}$ la sucesión de polinomios ortonormales asociada al producto de Sobolev en $(\Poly, \langle \cdot,\cdot\rangle_S)$ y sea
\(\mathbf D_n\) la sección truncada definida en
Definición~\ref{def:seccion_truncada}.
Entonces, para todo \(n\ge0\), existe una constante
\[
    \kappa_n\in\C\setminus\{0\}
\]
tal que
\[
    \det(z\mathbf I_{n+1}-\mathbf D_n^t)
    =
    \kappa_n\,\phi_{n+1}(z).
\]
En particular,
\[
    \sigma(\mathbf D_n^t)
    =
    \{z\in\C:\phi_{n+1}(z)=0\},
\]
contando multiplicidades algebraicas.
\end{prop}

\begin{proof}
	Como la multiplicación por $z$ aumenta el grado en una unidad, $\mathbf{D}_n$ tiene estructura de Hessenberg superior: sus entradas satisfacen $d_{i,j}=0$ para $i>j+1$. Esto implica que para cualquier $z\in\mathbb{C}$, resolviendo la relación matricial, se tiene
	\[
		\det(z\mathbf{I}_{n+1}-\mathbf{D}_n^t)=\det(z\mathbf{I}_{n+1}-\mathbf{D}_n)
		=\frac{1}{\gamma_{n+1}}\phi_{n+1}(z),
	\]
	donde $\gamma_{n+1}$ es el coeficiente director de $\phi_{n+1}$. La última igualdad es la identidad característica estándar para matrices de Hessenberg asociadas a polinomios ortogonales (véase por ejemplo \cite{szego1975orthogonal}, \S3.1). Por tanto, $z$ es autovalor de $\mathbf{D}_n^t$ si y solo si $\phi_{n+1}(z)=0$.
\end{proof}

\section{Valores singulares de \texorpdfstring{$\mathbf D_n$}{Dn}}


\begin{defn}\label{def:sigma_kn}
	Sea \(\mathbf D_n\in\mathcal M_{n+1}(\C)\) la sección truncada definida en
	Definición~\ref{def:seccion_truncada}.
	Para $$n\ge0$$, denotamos por
	\[
		\sigma_{0,n}\ge\sigma_{1,n}\ge\cdots\ge\sigma_{n,n}\ge0
	\]
	los valores singulares de \(\mathbf D_n\), ordenados de forma decreciente.
	En particular,
	\[
		\sigma_{0,n}=\|\mathbf D_n\|_2.
	\]
\end{defn}

Nótese que indexamos desde $0$ para que el índice $k$ de $\sigma_{k,n}$ coincida con el índice $k$ de $Q_k(\D)$, lo que facilitará la comparación entre ambos objetos.

\begin{lem}
\label{lem:minmax_sigma_Dn}
Sea \(n\ge0\). Para todo \(0\le k\le n\),
\[
    \sigma_{k,n}
    =
    \min_{\substack{S\subset\C^{n+1}\\ \dim S=n+1-k}}
    \max_{\substack{x\in S\\ x\neq0}}
    \frac{\|\mathbf D_nx\|_2}{\|x\|_2}.
\]
Equivalentemente,
\[
    \sigma_{k,n}
    =
    \max_{\substack{S\subset\C^{n+1}\\ \dim S=k+1}}
    \min_{\substack{x\in S\\ x\neq0}}
    \frac{\|\mathbf D_nx\|_2}{\|x\|_2}.
\]
\end{lem}

A continuación veremos la relación entre la fórmula min--max de Courant--Fischer y el índice $Q:k$.

\begin{prop}\label{prop:sigma_qk}
	Sea \(\mathbf D_n\) la sección truncada de \(\D\) en
	\(\mathcal P_n=\operatorname{span}\{\phi_0,\dots,\phi_n\}\).
	Entonces, para todo \(n\ge0\) y todo \(0\le k\le n\),
	\[
		\sigma_{k,n}\le Q_k(\D).
	\]
	En consecuencia,
	\[
		\lim_{n\to\infty}\sigma_{k,n}\le Q_k(\D).
	\]
\end{prop}

\begin{proof}
	Fijemos un subespacio admisible $V\subset\Poly$ con
	\[
		\operatorname{codim}_{\Poly}(V)<k+1.
	\]
	Para cada $n$, consideramos
	\[
		V_n:=V\cap\mathcal P_n.
	\]
	Al pasar a coordenadas en la base $\{\phi_0,\dots,\phi_n\}$ obtenemos un subespacio $S_n\subset\mathbb C^{n+1}$ con codimensión menor que $k+1$, o equivalentemente con dimensión al menos $n+1-k$. Por Courant--Fischer,
	\[
		\sigma_{k,n}
		\le
		\max_{\substack{x\in S_n\\x\neq0}}\frac{\|\mathbf D_nx\|_2}{\|x\|_2}.
	\]
	Si $x$ corresponde al polinomio $p\in V_n$, entonces
	\[
		\|\mathbf D_nx\|_2=\|P_n(\D p)\|_S\le \|\D p\|_S\le \|\D|_V\|_S\,\|p\|_S=\|\D|_V\|_S\,\|x\|_2.
	\]
	Por tanto,
	\[
		\sigma_{k,n}\le \|\D|_V\|_S.
	\]
	Tomando ínfimo sobre todos los subespacios admisibles $V$ con codimensión $<k+1$, obtenemos
	\[
		\sigma_{k,n}\le Q_k(\D).
	\]
	La afirmación sobre el $\lim$ es inmediata.
\end{proof}

\section{Conexión con el teorema de estabilización}

Ya podemos traducir el Teorema~\ref{thm:estabilizacion} al lenguaje matricial.

\begin{thm}\label{thm:umbral_singular_indice_m_enunciado}
	Sea
	\[
		\langle p,q\rangle_S
		=
		\int_{\mathbb S_R}p(z)\overline{q(z)}\,d\mathbf m_R(z)
		+
		\sum_{c\in I_{\mathrm{out}}} p'(c)\overline{q'(c)},
		\qquad p,q\in\Poly,
	\]
	y supongamos que $I_{\mathrm{in}}=\varnothing$ de la Definición \ref{def:class_points_new}. Sea $m=|I_{\mathrm{out}}|$. Entonces
	\[
		\lim_{n\to\infty}\sigma_{j,n}=+\infty,
		\qquad 0\le j\le m-1,
	\]
	y además
	\[
		\lim_{n\to\infty}\sigma_{m,n}\le R.
	\]
\end{thm}

\begin{proof}
	La prueba sigue el mismo principio de construcción que el Teorema~\ref{thm:estabilizacion}, pero ahora trasladado a las compresiones finitas.

	\smallskip
	\noindent\emph{Paso 1: explosión para $j<m$.}
	Basta probarla para $j=m-1$ por monotonía en el índice. En el Lema~\ref{lem:paquetes_diagonales} del capítulo anterior se construyeron, para cada $n$, polinomios
	\[
		u_{c,n}\in\Poly,
		\qquad c\in I_{\mathrm{out}},
	\]
	que satisfacen
	\[
		\|u_{c,n}\|_S\xrightarrow[n\to\infty]{}0,
		\qquad
		\psi_d(u_{c,n})=\delta_{dc}.
	\]
	Sea
	\[
		F_n:=\operatorname{span}\{u_{c,n}:c\in I_{\mathrm{out}}\}.
	\]
	Por la diagonalidad anterior, estos vectores son linealmente independientes, luego $\dim(F_n)=m$. Además, si
	\[
		p=\sum_{c\in I_{\mathrm{out}}}\beta_c u_{c,n}\in F_n,
		\qquad p\neq0,
	\]
	y definimos
	\[
		\psi^{(n)}:=\sum_{c\in I_{\mathrm{out}}}\overline{\beta_c}\,\psi_c,
	\]
	entonces, exactamente como en la prueba de la Proposición~\ref{prop:umbral_fase2}, se obtiene
	\[
		|\psi^{(n)}(p)|=\sum_{c\in I_{\mathrm{out}}}|\beta_c|^2
		\qquad\text{y}\qquad
		|\psi^{(n)}(p)|\le \|\D p\|_S.
	\]
	Si normalizamos $\|\beta\|_{\ell^2}=1$, la primera identidad da $|\psi^{(n)}(p)|=1$, mientras que
	\[
		\|p\|_S
		\le
		\sum_{c\in I_{\mathrm{out}}}|\beta_c|\,\|u_{c,n}\|_S
		\le
		\sqrt m\max_{c\in I_{\mathrm{out}}}\|u_{c,n}\|_S
		\xrightarrow[n\to\infty]{}0.
	\]
	Por tanto,
	\[
		\inf_{0\neq p\in F_n}\frac{\|\D p\|_S}{\|p\|_S}\xrightarrow[n\to\infty]{}+\infty.
	\]
	Fijado ahora $L>0$, elegimos $n(L)$ tan grande que
	\[
		\inf_{0\neq p\in F_{n(L)}}\frac{\|\D p\|_S}{\|p\|_S}>L.
	\]
	Sea $F:=F_{n(L)}$ y sea $N_L$ un entero tal que $F\subset \mathcal P_{N_L-1}$. Entonces, para todo $N\ge N_L$, se tiene $F\subset \mathcal P_{N-1}$ y por tanto $\D p\in\mathcal P_N$ para todo $p\in F$. En consecuencia, sobre $F$ la compresión coincide con $\D$, y la forma max--min de Courant--Fischer da
	\[
		\sigma_{m-1,N}\ge \inf_{0\neq p\in F}\frac{\|\D p\|_S}{\|p\|_S}>L,
		\qquad N\ge N_L.
	\]
	Como $L$ es arbitrario, concluimos que
	\[
		\sigma_{m-1,N}\xrightarrow[N\to\infty]{}+\infty.
	\]
	La monotonía de los valores singulares respecto del índice implica entonces
	\[
		\sigma_{j,n}\to+\infty,
		\qquad 0\le j<m.
	\]

	\smallskip
	\noindent\emph{Paso 2: cota en el índice $m$.}
	Sea
	\[
		V_{\mathrm{out}}:=\bigcap_{c\in I_{\mathrm{out}}}\ker(\psi_c).
	\]
	Por el Capítulo~\ref{sec:estabilizacion}, este subespacio tiene codimensión exacta $m$ y satisface
	\[
		\|\D p\|_S\le R\|p\|_S,
		\qquad p\in V_{\mathrm{out}}.
	\]
	Al intersectarlo con $\mathcal P_n$ obtenemos un subespacio de dimensión al menos $n+1-m$ donde la compresión $\mathbf D_n$ sigue estando acotada por $R$. De nuevo por Courant--Fischer,
	\[
		\sigma_{m,n}\le R,
	\]
	para todo $n$ suficientemente grande, y por tanto
	\[
		\lim_{n\to\infty}\sigma_{m,n}\le R.
	\]
	Esto completa la prueba.
\end{proof}

\begin{cor}\label{cor:estabilizacion_singular_umbral}
	Bajo las hipótesis del Teorema~\ref{thm:estabilizacion} se tiene:
	\begin{enumerate}
		 \item Si \(I_{\mathrm{in}}=\varnothing\), entonces
    \[
        \sigma_{j,n}\xrightarrow[n\to\infty]{}+\infty,
        \qquad 0\le j<m,
    \]
    y
    \[
        \limsup_{n\to\infty}\sigma_{m,n}\le R.
    \]

    \item En general,
    \[
        \limsup_{n\to\infty}\sigma_{m,n}
        \le
        Q_m(\D)
        <
        \infty.
    \]
	\end{enumerate}
\end{cor}

\begin{rmk}[Lectura matricial del umbral]
	El Corolario~\ref{cor:estabilizacion_singular_umbral} muestra que el índice
	\(m=|I_{\mathrm{out}}|\) es el primer índice para el que se obtiene una cota
	singular uniforme a partir del teorema de estabilización. En ausencia de
	átomos interiores, los valores singulares de índices \(j<m\) divergen, mientras
	que el índice \(m\) queda controlado por \(R\). En el caso mixto, el mismo índice
	queda controlado por \(Q_m(\D)\).
\end{rmk}

